\section{Introduction}\label{sec:introduction}

Write $[n]=\{1,\ldots,n\}$ for a positive integer $n$.
For a finite set $A\subset\N=\{1,2,\ldots\}$, write
\[
 \FS(A)=\left\{\sum_{a\in B}a:\varnothing\ne B\subseteq A\right\},
 \qquad
 \FP(A)=\left\{\prod_{a\in B}a:\varnothing\ne B\subseteq A\right\}.
\]
Thus each element of $A$ may occur at most once in an individual sum
or product, and singleton subsets are included. We prove the following.

\begin{theorem}\label{thm:main}
Let $r,m\ge1$ be integers, and fix real numbers $R\ge2$ and $D\ge1$.
For every coloring $\chi:\N\to[r]$,
there are distinct positive integers $a_1<\cdots<a_m$ and a color
$c\in[r]$ such that
\[
 \chi\left(\sum_{j\in J}a_j\right)
 =\chi\left(\prod_{j\in J}a_j\right)=c
 \quad\text{for every }\varnothing\ne J\subseteq[m].
\]
Equivalently, $\FS(A)\cup\FP(A)$ is monochromatic for a set
$A\subset\N$ of cardinality $m$.
The elements can additionally be chosen to satisfy
\begin{equation}\label{eq:separated-elements}
 a_1>R,\qquad
 a_d>R\left(\sum_{k<d}a_k+\prod_{k<d}a_k\right)^D
 \quad(2\le d\le m).
\end{equation}
\end{theorem}

Theorem~\ref{thm:main} resolves Hindman's finite sums and products
conjecture positively. The same assertion for colorings of
$\N_0=\{0,1,2,\ldots\}$ follows by restriction to $\N$.
The separation in \eqref{eq:separated-elements} also eliminates every
collision between expressions except the shared singleton values.

\begin{corollary}\label{cor:distinct-expressions}
For every finite coloring of $\N$ and every $m\ge1$, there is an
$m$-element set $A\subset\N$ for which $\FS(A)\cup\FP(A)$ is
monochromatic and
\[
 |\FS(A)|=|\FP(A)|=2^m-1,\qquad \FS(A)\cap\FP(A)=A.
\]
In particular, the union has $2(2^m-1)-m$ distinct elements.
\end{corollary}

The proof of the corollary is elementary and appears at the end of
Section~\ref{sec:framework}. Both conclusions concern prescribed
finite sets; the separation does not assert the existence of an
infinite simultaneous sequence.

\subsection{History and significance}

The separate additive and multiplicative problems have classical answers.
Schur's theorem guarantees a monochromatic triple $\{x,y,x+y\}$ in every
finite coloring of $\N$ \cite{Schur}.
The finite sums theorem of Folkman--Rado--Sanders extends this to the
nonempty subset sums of arbitrarily large finite sets; Sanders's original
argument appears in \cite[Theorem~2 and Corollary~1.1]{Sanders}.
Hindman proved the infinite finite sums theorem \cite{Hindman}.
Applying an additive theorem to the coloring $n\mapsto\chi(2^n)$ gives
its multiplicative counterpart. The simultaneous question is different:
Hindman constructed a finite coloring of $\N$ with no infinite set whose terms,
pairwise sums and pairwise products all have one color
\cite{Hindman1980}.
His finite conjecture asks for arbitrarily large finite sets instead
\cite{Hindman1979}; see also
\cite[Question~17.18]{HindmanStrauss}.

For two variables, computer-assisted arguments of Graham and Hindman
already established the existence of a monochromatic quartet
$\{x,y,x+y,xy\}$ in every two-coloring of $\N$; see the historical
account in \cite[p.~82]{HindmanPhulara}.
Bowen gave a proof in which $x$ and $y$ are distinct and exceed any
prescribed bound, and more generally controlled the individual variables, their prefix
products and their total sum \cite[Theorem~1.1]{BowenTwoColors}.
For arbitrary finite colorings, Moreira proved the existence of
monochromatic triples $\{x,x+y,xy\}$ as part of a wider theorem on
prefix products and polynomial shifts \cite[Theorem~1.4 and
Corollary~1.5]{Moreira}.
Alweiss gave a polynomial proof of the mixed triple and related
polynomial-shift patterns \cite[Theorem~1.2]{AlweissPolynomials}.
More recently, Alweiss, Bowen and Sabok proved that every two-coloring
contains $\{x,y,xy,x+iy:1\le i\le k\}$ in one color for every prescribed
$k$ \cite[Theorem~1.2]{AlweissBowenSabok}.

Results over fields developed in parallel. Green and Sanders proved
the four-term result in sufficiently large prime fields
\cite[Theorem~1.2]{GreenSanders}, while Bergelson and Moreira developed
affine ergodic methods for mixed additive--multiplicative patterns
\cite{BergelsonMoreiraField,BergelsonMoreiraAffine}.
Bowen and Sabok established $\{x,y,x+y,xy\}$ over $\Q$
\cite[Theorem~1.1]{BowenSabok}, and Alweiss proved the full finite sums
and products theorem over $\Q$ \cite[Theorem~1.3]{Alweiss}.
Richter obtained integer density theorems for the mixed
pair $\{x+Q(y),xy\}$ using Fourier analysis, ergodic theory and
logarithmic averages over products of primes \cite[Theorems~1.5 and
1.7]{Richter}.

The passage from rational to integer configurations requires more than
clearing denominators. A common dilation scales a sum by one power of
the dilation and a product of $j$ entries by its $j$th power; an
arbitrary coloring need not respect these different changes.
The ordered-block selection in our proof follows the Ramsey and
finite-sums argument in \cite[Section~5]{Alweiss}.
Our Alignment argument also adapts the compensating scale updates in
\cite[Section~4]{Alweiss}, which fix the current target product while
retaining the requirements secured at earlier targets.
The additional task is to control the additive shifts arithmetically
while preserving all the required multiplicative colors.

\subsection{Proof strategy}

The finite combinatorial part of the argument uses ordered blocks of
widely separated integer variables. A Ramsey argument and the finite
sums theorem select block products whose nonempty products already
have one color. The remaining task is to ensure that their sums have
that same color. We do this by constructing bounded predictions for
the color indicators and then arranging that those predictions remain
positive at the required sums.

The proof separates these tasks into a Prediction Principle and an
Alignment Principle. Prediction replaces the color functions inside
the relevant counts by piecewise nilsequences (Lipschitz functions
evaluated along nilpotent orbits, separately on intervals and residue
classes), while retaining the multiplicative color masks.
It also ensures that a color occurring at
a center is unlikely to have a small predicted value there. Alignment
selects from a fixed finite list of rational scale vectors so that a
positive prediction at every center remains positive after all the
required additive shifts. Its success probability is bounded below
independently of the complexity of the nilsequence models. This
uniformity permits the final choice of small calibration errors.

Three constructions provide the analytic and arithmetic inputs.
\begin{enumerate}
\item \emph{Weights for the distribution of a block product.}
Each chosen element is a product of several raw variables, so its
distribution differs from that of a single variable at the same largest
scale. We encode divisibility by the earlier variables in nonnegative
weights and attach the same weight to every sum ending at that block.
Prime substitutions and weighted Cauchy--Schwarz remove the product-color
conditions and reduce the counting error to one-variable additive tests.

\item \emph{Predictions at two additive scales.}
The counting estimate tests short additive shifts, whereas Alignment
needs one model valid across a larger interval. We construct nilsequence
models on nested intervals and use Ramsey selection to make their
projections close. This supplies a common model for both requirements,
with nilpotence step depending only on the requested configuration size.

\item \emph{Realizing shifts through nilsequence states.}
Adding an earlier block product may require a noninteger change in the
largest-scale variable of that block. We first make an integer residue correction and then
realize the remaining displacement as a transformation of the
nilsequence model. Two separate arguments justify this transformation:
comparison of averages on arithmetic progressions, and a lifting theorem
that makes transformations for different models act on one common state.
These transformations generate a nilpotent group whose step is independent
of model complexity. Finite polynomial recurrence in this group gives
the required uniform alignment probability.
\end{enumerate}

The rough-step comparison and the passage from marginal cube symmetries
to joint lifts are stated separately in
Propositions~\ref{prop:rough-step} and~\ref{prop:face-lift}. Their
hypotheses make explicit the joint polynomial dependence and the
cube information needed to control otherwise unavailable shifts.
The proof uses the inverse theorem of Green, Tao and Ziegler
\cite{GTZ,GTZErratum}, the concatenation theorem of Tao and Ziegler
\cite{TaoZiegler}, quantitative polynomial equidistribution
\cite{GreenTao,GreenTaoErratum}, and nilpotent polynomial recurrence
\cite{ZorinKranich}. We state the required forms when they are used.
The intervening transference, progression-comparison and lifting
arguments are proved below.

Section~\ref{sec:framework} gives the two principles and derives
Theorem~\ref{thm:main} from them, including integrality and separation.
Sections~\ref{sec:arithmetic}--\ref{sec:prediction} develop Prediction;
Sections~\ref{sec:rough-progressions}--\ref{sec:alignment} develop Alignment.
The missing-corner Lemma~\ref{lem:cube-corner}, proved independently in
Section~\ref{sec:cube-limits}, is also used in Prediction. All parameters
are qualitative: their finiteness and order of choice are essential,
but no numerical bound for the smallest configuration is claimed.

\begin{figure}[htbp]
\centering
\begin{tikzpicture}[
 node distance=8mm and 7mm,
 box/.style={draw=black!45,rounded corners=2pt,align=center,
 text width=5.65cm,minimum height=12mm,inner sep=5pt,font=\small},
 result/.style={box,text width=9.8cm},
 arr/.style={-{Stealth[length=2mm]},semithick,draw=black!60}]
 \node[box] (weights) {Arithmetic scales and divisor weights\\
   Section~\ref{sec:arithmetic}};
 \node[box,right=of weights] (geometry) {Rough progressions and cube lifts\\
   Sections~\ref{sec:rough-progressions}--\ref{sec:cube-limits}};
 \node[box,below=of weights] (prediction) {Weighted counts and bounded models\\
   $\Downarrow$\\Prediction, Sections~\ref{sec:correlation}--\ref{sec:prediction}};
 \node[box,below=of geometry] (alignment) {Finite polynomial recurrence\\
   $\Downarrow$\\Alignment, Section~\ref{sec:alignment}};
 \path (prediction.south) -- (alignment.south)
   node[midway,below=13mm,result] (deduction)
   {Product-chain selection and a positive weighted count\\
    Theorem~\ref{thm:main} and Corollary~\ref{cor:distinct-expressions}};
 \draw[arr] (weights) -- (prediction);
 \draw[arr] (geometry) -- (alignment);
 \draw[arr] (geometry.south west) -- (prediction.north east);
 \draw[arr] (prediction.south) -- (deduction.north west);
 \draw[arr] (alignment.south) -- (deduction.north east);
\end{tikzpicture}
\caption{The main analytic routes meet in the combinatorial deduction.
The diagonal arrow records the shared missing-corner reconstruction
from Lemma~\ref{lem:cube-corner}.
Alignment fixes a positive mass before Prediction fixes model complexity;
this order permits the counting error to be made small enough.}
\label{fig:proof-map}
\end{figure}
